\ulohahead{společná informatika}{Objektové programování a návrh (C\#)}
\begin{zadani}
\begin{enumerate}[label=(\alph*)]
  \item \bod{1} Vysvětlete abstrakci, zapouzdření a polymorfismus a související konstrukty (třída, rozhraní, dědičnost, viditelnost). Rozdíl mezi statickým a dynamickým typováním a mezi virtuální a nevirtuální metodou.
  \item \bod{2} Navrhněte v C\# malou hierarchii geometrických tvaru s rozhraním a polymorfním výpočtem obsahu; ukažte ošetření chyby výjimkou.
  \item \bod{3} Jak je implementován dynamický polymorfismus (tabulka virtuálních metod)? K čemu jsou generické typy a jaký je rozdíl mezi hodnotovými a referenčními typy v C\#?
\end{enumerate}
\end{zadani}

\begin{reseni}
\textbf{(a)} \emph{Abstrakce} skrývá detaily za rozhraním; \emph{zapouzdření} drží data + operace pohromadě a omezuje přístup (viditelnost \texttt{private}/\texttt{public}/\texttt{protected}); \emph{polymorfismus} = jednotné volání nad různými typy. \emph{Třída} definuje typ, \emph{rozhraní} (interface) jen kontrakt, \emph{dědičnost} přebírá a rozšiřuje. \emph{Statické} typování kontroluje typy při překladu, \emph{dynamické} za běhu. \emph{Virtuální} metoda se vybírá podle skutečného typu objektu za běhu (override), \emph{nevirtuální} podle deklarovaného typu.

\textbf{(b)}
\begin{lstlisting}
public interface IShape { double Area(); }

public sealed class Circle : IShape {
    private readonly double r;
    public Circle(double r) {
        if (r < 0) throw new ArgumentException("zaporny polomer");
        this.r = r;
    }
    public double Area() => Math.PI * r * r;
}

public sealed class Rectangle : IShape {
    private readonly double w, h;
    public Rectangle(double w, double h) { this.w = w; this.h = h; }
    public double Area() => w * h;
}

double TotalArea(IEnumerable<IShape> shapes) {
    double sum = 0;
    foreach (var s in shapes) sum += s.Area(); // polymorfni volani
    return sum;
}
\end{lstlisting}

\textbf{(c)} \emph{Dynamický polymorfismus} se realizuje \emph{tabulkou virtuálních metod} (vtable): každý objekt drží ukazatel na tabulku své třídy, volání virtuální metody = skok přes položku tabulky (vybere se override potomka). \emph{Generické typy} umožní psát kód nezávislý na konkrétním typu (\texttt{List<T>}) bez ztráty typové kontroly a bez boxingu u hodnotových typu. \emph{Hodnotové} typy (\texttt{struct}, \texttt{int}) se kopírují podle hodnoty (typicky na zásobníku), \emph{referenční} (\texttt{class}) žijí na haldě a předává se reference; spravuje je garbage collector.
\end{reseni}

\kalibrace{Programovací jazyky / OO návrh je v moderním formátu vzácnější, ale když padne, bývá implementační. Definice (a) + funkční kostra hierarchie (b) tě udrží nad nulou; jazyk je C\# (volba majitele).}
\past{Rozhraní není abstraktní třída (nemá stav/implementaci, lze jich implementovat víc). Virtuální dispatch jde podle \emph{skutečného} typu - jinak by polymorfismus nefungoval.}
